\subsection{\_time}\label{time}
The \verb|_time| module provides a lightweight run-time measurement
(profiling). It is based on the cycle counter of the ARM Cortex-M debug unit
(\verb|DWT->CYCCNT|), which increments on every processor clock and thus
allows a resolution in the nanosecond range. Up to \verb|TIME_DEV_CNT|
independent, named measurement channels (\emph{handles}) are managed; each
channel holds a pool of \verb|TIME_MEAS_CNT| individual measurements.

\subsubsection*{Initialisation}
\verb|time_init()| must be called before first use. It activates the trace
system in the debug-core register and starts the cycle counter (Listing
\ref{time_init_regs}). From the clock frequency obtained via
\verb|HAL_RCC_GetHCLKFreq()| the conversion factor \verb|ccnt2ns|
($= 10^9 / f_{\mathrm{HCLK}}$) is computed, with which the macro \verb|NOW|
converts the current counter value into nanoseconds.

\begin{listing}
\begin{minted}{C}
CoreDebug->DEMCR |= CoreDebug_DEMCR_TRCENA_Msk; // enable trace
DWT->CTRL        |= DWT_CTRL_CYCCNTENA_Msk;     // enable cycle counter
\end{minted}
\caption{Enabling the DWT cycle counter in \texttt{time\_init()}}
\label{time_init_regs}
\end{listing}

\begin{listing}
\begin{minted}{C}
#define TIME_MEAS_CNT             10  // measurements per handle
#define TIME_MEAS_CHAR_PER_LINE   20  // max. length of the label string
#define TIME_MEAS_TIME_FIELD_LEN  12
#define TIME_DEV_CNT               6  // number of concurrent handles
\end{minted}
\caption{Configuration constants from \texttt{\_time.h}}
\label{time_defines}
\end{listing}

\subsubsection*{Modes}
The behaviour of a handle is controlled by a bit field of \verb|mode_e|
(Listing \ref{mode_e}), set with \verb|time_set_mode()|:
\begin{description}
\item[\texttt{ONESHOT}] Exactly one pool pass is recorded; afterwards
  \verb|doLoop| is set to \verb|false| and recording stops (queryable via
  \verb|time_doLoop_get()|).
\item[\texttt{MAX\_HOLD}] Keeps the lowest \verb|max| pool entries (cf.\
  \verb|time_set_max()|); only the entries above are overwritten cyclically.
\item[\texttt{SSSC}] \emph{Store subslot, slot, cycle}: instead of the
  supplied label string, the cycle position is stored in the data field via
  \verb|cycle_string()| (see Section \ref{cycle}).
\end{description}

\begin{listing}
\begin{minted}{C}
typedef enum {
    ONESHOT  = 1,
    MAX_HOLD = 2,
    SSSC     = 4, // store subslot, slot, cycle data in data field
} mode_e;
\end{minted}
\caption{Operating modes \texttt{mode\_e}}
\label{mode_e}
\end{listing}

\subsubsection*{Data structures}
A single measurement is held in \verb|time_meas_t| (Listing
\ref{time_meas_t}). It stores both the raw tick values (\verb|tick_start|,
\verb|tick_stop|) and the times derived from them in nanoseconds. Two stop
instants are distinguished: an intermediate stop \verb|stop_su_ns|
(\emph{setup}, set by \verb|time_stop_su()|) and the final stop
\verb|stop_tx_ns|. The field \verb|line| holds the label string or the SSSC
cycle information.

\begin{listing}
\begin{minted}{C}
typedef struct time_meas_s {
    uint32_t tick_start;
    uint32_t tick_stop;
    uint32_t tick_duration;
    int64_t start_ns;
    int64_t stop_su_ns;
    int64_t stop_tx_ns;
    int64_t duration_su_ns;
    int64_t duration_tx_ns;
    int64_t baud;
    uint8_t count;
    char line[TIME_MEAS_CHAR_PER_LINE + 1];
} time_meas_t;
\end{minted}
\caption{Definition of the \texttt{time\_meas\_t} structure}
\label{time_meas_t}
\end{listing}

All measurements of a handle are gathered in \verb|time_single_t|
(Listing \ref{time_single_t}). Besides the measurement pool
\verb|measurement[]| it holds the current ring-buffer index \verb|idx|, the
lower hold bound \verb|max| (for \verb|MAX_HOLD|), aggregate statistics
(\verb|max_baud|, \verb|max_cnt|) and the channel name.

\begin{listing}
\begin{minted}{C}
typedef struct time_single_s {
    time_meas_t measurement[TIME_MEAS_CNT];
    int8_t idx;
    uint8_t max;
    int64_t init_ns;
    int64_t last_start_ns;
    int64_t first_ns;
    int64_t last_ns; // set once the maximum baud rate was reached
    float max_baud;
    int32_t max_cnt;
    uint8_t mode;
    uint8_t name[LINE_CHAR];
} time_single_t;
\end{minted}
\caption{Definition of the \texttt{time\_single\_t} structure}
\label{time_single_t}
\end{listing}

The enclosing structure \verb|timem_t| (Listing \ref{timem_t}) is allocated in
\verb|_time.c| as the single static variable \verb|_time|. It manages the
array of all handles, the allocation bitmask \verb|used|, the clock frequency
\verb|clk_Hz| and the conversion factor \verb|ccnt2ns|.

\begin{listing}
\begin{minted}{C}
typedef struct timem_s {
    time_single_t time[TIME_DEV_CNT];
    uint8_t used;       // allocation bitmask of the handles
    uint32_t clk_Hz;
    float ccnt2ns;      // tick -> ns
    int64_t init_ns;
    bool init;
    bool doLoop;        // true while recording (ONESHOT)
} timem_t;
\end{minted}
\caption{Definition of the \texttt{timem\_t} structure}
\label{timem_t}
\end{listing}

\subsubsection*{API}
\begin{description}
\item[\texttt{time\_init()}] Initialises the module: enables the DWT cycle
  counter (Listing \ref{time_init_regs}), computes \verb|ccnt2ns| from the
  HCLK frequency and resets all handles. Must be called before any other
  function.
\item[\texttt{time\_new(name)}] Reserves a free handle, stores the name and
  returns its \verb|time_handle_t|, or $-1$ if all \verb|TIME_DEV_CNT| handles
  are taken.
\item[\texttt{time\_set\_mode(hdl, mode)} / \texttt{time\_get\_mode(hdl)}]
  Sets or reads the operating-mode bit field (\verb|mode_e|) of the handle.
\item[\texttt{time\_set\_max(hdl, max)} / \texttt{time\_get\_max(hdl)}]
  Defines the number of retained pool entries ($0 \le \mathtt{max} <
  \mathtt{TIME\_MEAS\_CNT}-1$); the remaining slots are overwritten
  cyclically.
\item[\texttt{time\_reset(hdl)}] Clears the handle's measurements but keeps the
  initialisation instant \verb|init_ns|.
\item[\texttt{time\_start(hdl, count, ptr, cycle)}] Starts a measurement:
  stores start tick and time, the quantity \verb|count| (basis of the baud
  computation) and the label \verb|ptr| -- or, with \verb|SSSC| set, the cycle
  position from \verb|cycle|.
\item[\texttt{time\_stop\_su(hdl)}] Records the intermediate stop and computes
  \verb|duration_su_ns| (time from start to \emph{setup}).
\item[\texttt{time\_stop(hdl, ptr)}] Records the final stop, computes
  \verb|duration_tx_ns|, \verb|baud| and the tick duration and advances the
  ring-buffer index. Optionally \verb|ptr| updates the label.
\item[\texttt{time\_auto(hdl, count, ptr, cycle)}] Convenience:
  \verb|time_start()| immediately followed by \verb|time_stop()|.
\item[\texttt{time\_now\_ns()}] Returns the current counter value in
  nanoseconds (\verb|NOW|).
\item[\texttt{time\_doLoop\_get()}] Returns \verb|doLoop|; in \verb|ONESHOT|
  mode it becomes \verb|false| after one full pass.
\item[\texttt{time\_print(hdl, titel, python, timing)}] Prints the pool as a
  table. See the next section.
\end{description}

\subsubsection*{Measurement flow and output}
{\sloppy\emergencystretch=3em
A typical measurement consists of \verb|time_start()| and a matching
\verb|time_stop()| (or the combined \verb|time_auto()|); optionally an
intermediate time can be taken with \verb|time_stop_su()|. From the duration
and the supplied quantity the baud rate is computed:
\[
  \mathtt{baud} = \frac{\mathtt{count} \cdot 8 \cdot 10^{9}}
                       {\mathtt{duration\_tx\_ns}} \quad [\mathrm{bit/s}]
\]
The measurements are stored in a ring buffer: entries below \verb|max| are
retained, the rest are overwritten cyclically. In \verb|ONESHOT| mode
\verb|doLoop| is set to \verb|false| after the first full pass.

\verb|time_print()| prints the pool as a table. With \verb|python == true| the
serial interface is switched to \verb|RAW| mode for the output and the result
is formatted as a Python dictionary so that it can be read directly by the
analysis scripts. The parameters \verb|timing| and the \verb|SSSC| flag select
which columns (plain times or cycle/slot/sub-slot) are printed.
\par}
